% ============================================================================
% Abstract --- Technical Reference
% Format: 450--900 characters, three layers (what / what it says / why it matters)
% ============================================================================

\noindent This document is a practitioner-oriented technical reference
and migration playbook for universities and research-performing
institutions seeking sovereignty over their digital infrastructure. It develops a
five-dimension sovereignty grid --- Data, Technical, Financial,
Operational, Institutional --- with twenty measurable indicators on a
four-point scale, and introduces the \emph{Trusted Activity Context}
(TAC) as the unit of reasoning that binds subject, action, resource
and context across layers. It then establishes a vendor-neutral
reference architecture organised in nine layers across three planes,
demonstrated on open foundations and admitting proprietary components
on the condition that they honour the published interface contracts.
A ten-wave migration playbook supports gradual evolution from any
existing brownfield estate, with three trajectories per wave
(greenfield, brownfield single-suite-centric, brownfield hybrid) and
explicit reversibility criteria; a structured conformance suite
provides the evidence base for closure. The document closes with a
regulatory mapping that aligns the architecture and the playbook to
the institution's applicable data-protection, network-and-information-security,
trust-services and AI-governance regime, together with its
cross-border data-transfer regime --- with GDPR, NIS2, the EU AI Act,
eIDAS~2 and the Schrems~II line of case law as a worked instantiation
--- alongside the jurisdiction-neutral standards spine (NIST CSF~2.0,
ISO/IEC 27001/27002/27005, COBIT~2019). The architecture, the
playbook and the mapping are designed to generalise across
institutions and jurisdictions and are intended as the shared technical
baseline for the international university network anticipated in phase~2
of the SUIT initiative.
